\begin{proposition}\label{prop:decic57-exclusion}
There is no elliptic curve over a number field of degree $10$
with a point of order $57$.
\end{proposition}
\begin{proof}
Put $X=X_1(57)$. A hypothetical level structure gives an effective
rational trace divisor $D$ of degree $10$ on $X$, with noncuspidal
support. We allow repetitions in this trace. The curve has genus
$85$ and good reduction at $2$.

We first record the degree-independent inputs used below.
The exact winding computation accounts for every simple factor of
$J_X$ except the elliptic factor $57\mathrm a1$. On that factor the
Hecke eigenvalues at $2,5,7,11,13$ are respectively
\[
 -2,-3,-5,1,2,
\]
and every diamond acts trivially. Thus every diamond difference
has finite image on $J_X(\mathbb Q)$. We do not assume that
$J_X(\mathbb Q)$ has rank zero.
The global and residual cusp signatures are
\[
 1^{18}2^9 18^2,\qquad 1^{18}2^9 9^2 18.
\]
In particular every cusp of degree at most $4$ is globally intact.

Let $Y$ be the genus-$9$ quotient by the subgroup generated by
$4$ in $(\mathbb Z/57\mathbb Z)^*/\{\pm1\}$. We use its rational
$q$-expansion, or multiplicative-level, model. The map $X\to Y$
is defined over $\mathbb Q$ and has degree $9$; its displayed
$j$-function is the invariant of the quotient elliptic curve.
Write $\rho$ for its diamond involution induced by $5$, and put
$\alpha=\rho-1$. The image abelian fourfold has rank zero and
has $49$ and $441$ points over $\mathbb F_2$ and $\mathbb F_5$.
The proven finite Jacobian and its $\alpha$-image are
\[
 J_Y(\mathbb F_2)\simeq C_{105}\times C_{210},\qquad
 \alpha J_Y(\mathbb F_2)\simeq C_7\times C_7.
\]
For the multiplicative-level correspondence,
$T_5^\mu=\rho\operatorname{Frob}_5+5\operatorname{Frob}_5^{-1}$.
The operator $T_5^\mu-\rho-5$ annihilates the rational finite
$\alpha$-image. Its kernel on the displayed finite image is a
subgroup $K\simeq C_7$. These computations use the saturated
canonical model, its exact $j$-map of degree $160$, and a proven
class-group presentation. In the retained cusp order, the
$j$-widths, diamond permutation, and Galois permutation at $5$ are
\[
 (1,19,57,3,57,1),\quad (6,2,5,4,3,1),\quad
 (1,2,5,4,3,6).
\]
All classes in this test are normalized at the first rational cusp.

The Hasse bound excludes noncuspidal residual places of degree
at most $5$. Hence the noncuspidal part of $\overline D$, if
nonempty, is a single place $P$ of degree $d\in\{6,7,8,9,10\}$,
with multiplicity one. Its cuspidal complement has degree at
most $4$, so lifts to a rational cuspidal divisor. Subtracting its
normalized $\alpha$-class shows that the normed image of $P$
on $Y$ must pass the $K$ test.

The degree-independent tests in degrees $6,7,8,9$ exclude degrees
$7,9$ and leave only the two supersingular source models in degree
$6$ and one ordinary Galois orbit in degree $8$. Their exclusion
is recalled below. For degree $10$ we perform a fresh complete
census of all $2053$ twists over $\mathbb F_{1024}$. Exactly $60$
have a point of order $57$, giving six closed $j$-polynomials,
all of degree $10$. Since $57$ is squarefree, this condition is
preserved by isogeny. Thus the same census contains all possible
$j$-values on the dual quotient.
There are $120$ closed places of degree $10$ on $Y$. Together
with the places of degrees $1,2,5$, these give a containing set
of $48$ normed images; the $K$ test leaves six places over the
three polynomials
\[
 \begin{gathered}
 x^{10}+x^7+1,\\
 x^{10}+x^9+x^7+x^3+1,\\
 x^{10}+x^7+x^6+x^3+1.
 \end{gathered}
\]
The only entirely cuspidal residual divisors not globally intact
have a nonic cusp and a rational cusp. On $Y$ their normalized
classes are $9\alpha(c)+\alpha(r)$, where $c$ is either
width-$57$ cusp and $r$ either width-$1$ cusp. Both
$9\alpha(c)$ lie outside $K$, while every $\alpha(r)$ lies in
$K$. These reductions are excluded as well.

We next exclude the noncuspidal possibilities on the source.
Here we use the point-level correspondence and write
$\sigma=\langle5\rangle$. Its Eichler--Shimura formula is
\[
 T_5^{\rm pt}=\operatorname{Frob}_5+
                 5\sigma\operatorname{Frob}_5^{-1}.
\]
The rational class $[\sigma D-D]$ is torsion: its normalized
part has finite diamond-difference image, and its remaining
cusp difference is torsion by Manin--Drinfeld. The same is true
of the globally rational cuspidal complement. Reduction of
rational torsion is injective at the good odd prime $5$.
Consequently, putting $T=T_5^{\rm pt}P$, the divisor
\[
 B=(\sigma-1)(T_5^{\rm pt}-1-5\sigma)P
   =\sigma T-T+P+4\sigma P-5\sigma^2P
\]
must be principal.

In degrees $6,8$ the retained genuine local arcs and exact
five-isogenies give nonprincipal divisors $B$ of positive and
negative degrees $42,42,88$ for the three source representatives.
This calculation is independent of the degree of $D$; the only
use of that degree was the globally rational cuspidal complement,
which remains valid here.

For degree $10$ we match the physical source to the dual model
explicitly. Over all six physical $j$-orbits, the groups are
$C_{1026}$ in four cases and $C_3\times C_{342}$ in two.
Enumeration of every order-$57$ point gives $12$ diamond orbits
of cyclic subgroups. For each, we compute the invariant of the
quotient by its actual cyclic $57$-subgroup. Exactly six have
dual $j$-polynomial in the displayed containing set. This
matching is independently checked by enumeration of all rational
points and by composing actual $3$- and $19$-isogenies.
In particular, we do not identify the physical and dual
$j$-polynomials, or assume cyclicity of the full $57$-torsion.
The six required orbits occur among eight retained numerical
tests. In each test the exact five-kernel orbits have closed
target degrees $10,50$ and norm multiplicities $1,1$.
Every actual point image is retained, including its multiple by
$5$. The products of the kernel norms equal the full
$5$-division polynomial. Equal closed places are identified
before adding signed multiplicities; each resulting $B$ has
positive and negative degree $110$.

For completeness, the nonprincipality test uses
$R_i=H^0(X,iK_X)$, with
\[
 \dim R_1=85,\qquad \dim R_2=252,\qquad \dim R_4=588.
\]
All $3655$ and $31878$ product identities are verified through
$q^{1100}$, beyond the respective Sturm bounds $480,960$.
The modular units have pole degrees $37,39$, so the canonical
and $2K_X$ unit identities are global. For each new ordinary
source and target we deform $a_6$ to $a_6+t$, lift the
$57$-division equation and elliptic equation, and recover the
$2K_X$ arc through order $16$. Independent recovery of the
canonical arc gives rank $84$ for its unit equations and a
nonzero tangent. All $56$ retained arcs are therefore genuine.
Choose nonzero $F,H\in R_2$ vanishing on $B_+$. If
$f\in\mathcal L(B)$, then $G=fF$ belongs to $R_2$, vanishes on
$B_-$, and satisfies $GH\in FR_2$. In all eight tests, the
vanishing conditions on each sign have rank $110$,
multiplication by $F$ has rank $252$, and the joint map
\[
 R_2\longrightarrow
 H^0(B_-,2K_X|_{B_-})\oplus R_4/FR_2,
 \qquad G\longmapsto(G|_{B_-},GH\bmod FR_2)
\]
has rank $252$. Thus $\mathcal L(B)=0$, a contradiction.
The full calculation is replayed independently in Sage.

It remains to exclude globally intact cuspidal reduction.
We work on the full Jacobian in the multiplicative-level
$q$-expansion model, whose diamond matrix corresponds to physical
point multiplication by $5^{-1}$. For
$n\in\{2,5,7,11,13\}$ and $\ell\in\{5,7,11,13\}$ put
\[
 \Phi_{n,\ell}=(T_n-a_n)
                    (T_\ell-\langle\ell\rangle-\ell).
\]
The first factor has finite image on $J_X(\mathbb Q)$, and
the second kills rational torsion. Hence every $\Phi_{n,\ell}$
annihilates $J_X(\mathbb Q)$. Their pullbacks of differentials
span a space of dimension $82$ over $\mathbb F_2$.
This calculation takes place on the full integral canonical
lattice and uses no cotangent exactness for a quotient at $2$.

Let $C$ be the rational cuspidal divisor lifting $\overline D$.
Then all $\Phi_{n,\ell}[D-C]$ vanish. We test formal immersion
of their product on $X^{(10)}$ at $\overline C$.
The rational cusp expansion is retained through $q^{1100}$;
the genuine quadratic cusp arc comes from
$E(q^{19}),P=\zeta_3q$ and is verified through order $20$.
Diamonds give all $18$ rational and $9$ quadratic cusp places;
on the latter arc the ninth power is coefficientwise Frobenius.
At a quadratic place we use both $\mathbb F_2$-trace components.

If a rational cusp occurs, move one of maximum multiplicity
$n_0$ to the first cusp. The other rational multiplicities are
at most $\min(n_0,5)$ and the quadratic multiplicities at most
$5$. Complete enumeration gives, for $n_0=1,\ldots,10$,
\[
 598417,1165590,455634,124979,29835,
 6267,1122,162,17,1
\]
configurations. There are a further $\binom{13}{8}=1287$
divisors supported entirely at quadratic cusps. Every one of
these $2383311$ matrices has rank $10$; the enumeration counts
also agree with independent generating-function coefficients.
The product map is therefore a formal immersion at
$\overline C$. Since $D$ and $C$ have the same reduction and
the same image, formal immersion gives $D=C$. This contradicts
the noncuspidal support of $D$ and completes the proof.
\end{proof}
